\documentclass{article}

\usepackage[T1]{fontenc}

\usepackage{booktabs,graphicx,siunitx}

\begin{document}

\section{Friedman}

\begin{tabular}{ll}
\toprule
Condition & Mean rank \\
\midrule
score\_t1 & 1.23 \\
score\_t2 & 2.04 \\
score\_t3 & 2.73 \\
\bottomrule
\end{tabular}

\begin{tabular}{llll}
\toprule
$\chi$\textsuperscript{2} & df & p & Kendall's W [95\% CI] \\
\midrule
67.19 & 2 & $<$.001 & 0.56 [0.48, 1.00] \\
\bottomrule
\end{tabular}

\begin{tabular}{ll}
\toprule
Pair & padj \\
\midrule
score\_t1 - score\_t2 & $<$.001 \\
score\_t1 - score\_t3 & $<$.001 \\
score\_t2 - score\_t3 & $<$.001 \\
\bottomrule
\end{tabular}

\includegraphics{figures/01_friedman/figure_profile.png}

The nonparametric counterpart to repeated-measures ANOVA. The $\chi$\textsuperscript{2}/p tests whether ranks differ across conditions; Kendall's W is the effect size (0--1 agreement). Post-hoc shows which conditions differ. Your significance threshold ($\alpha$) is 0.05. N = 60 subjects.

A Friedman test gave $\chi$\textsuperscript{2}(2)=67.19, p $<$ .001, W=.56 [.48, 1.00].

Statistical basis: Friedman rank test (nonparametric counterpart to repeated-measures ANOVA) (Friedman, M. (1937). "The use of ranks to avoid the assumption of normality implicit in the analysis of variance." Journal of the American Statistical Association, 32(200), 675-701.); Kendall's W effect size (agreement among rankings) (Kendall, M. G. (1948). Rank Correlation Methods. Charles Griffin \& Company.); Nemenyi post-hoc pairwise comparisons (Nemenyi, P. B. (1963). Distribution-free multiple comparisons [Doctoral dissertation, Princeton University].).

\end{document}